\section{The witness count with the extended capacity}\label{adp:counts}

We apply Theorem~\ref{adp:plain} to the physical rows and witnesses
of \cite[Sections 8 and 19]{OAI26}. To distinguish two uses of $q$,
write $q_{\rm amp}$ for the mean prime amplitude; the moving-twist
length in this application is zero. Each row character has the fixed
ray twist $\nu$ and the varying factor $\chi_n(u)$.

\begin{proposition}[Actual witness count]\label{adp:witness-count}
Suppose $13/18\le\kappa\le3/4$ and
$\beta_*\le(1+\kappa)/2$. Fix a dynamic and amplitude bin of
retained sixth-power-free physical rows $q_u\asymp U$, with
\[
 1/50<\delta=2a-1\le2\beta_*-1\le\kappa,\qquad
 q_{\rm amp}\in[0,\delta/2],\qquad x=q_{\rm amp}/\delta .
\]
Assume the total available prime length at base $U$ exceeds $7/37$
by a fixed positive amount, with disjoint underlying prime windows
and the coefficients of \cite[Section 19.2]{OAI26}.
For $t\in[1,3/2]$, put
\[
 c=\frac1{3\kappa},\quad B_x=2-2cx,\quad
 D_x=B_x+1-x,\quad P_x=B_x(1-x),\quad \alpha=\frac56 .
\]
Then, for every prescribed $\varepsilon>0$, the strict capacity
decrements and the prime mesh can be chosen so that
\begin{equation}\label{adp:count}
 \#\mathcal B\ll U^{\max\{R_{\rm short}(t),\,L(t)\}+\varepsilon}
                    (1+T_1)^{A_\varepsilon},
\end{equation}
where
\[
 R_{\rm short}(t)=1-\delta+
             \frac{\delta P_x}{D_x}(3/2-t),\qquad
 L(t)=1-\delta+(\alpha-\delta)(t-1).
\]
The estimate uses the actual rowwise witnesses, profiles, cutoffs
and height condition of \cite[Proposition 8.3]{OAI26}.
All moment, selection, rounding and parameter-recovery costs are
included in $\varepsilon$ and the fixed polynomial height factor.
\end{proposition}

\begin{proof}
Subdivide by the witness presentation and its dyadic length pair,
as in \cite[Proposition 19.2]{OAI26}. The logarithmic number of
subdivisions is absorbed in the prescribed loss. For one subdivision
write $r,m$ for the inverse and plain witness lengths. The detector
gives both inequalities
\begin{equation}\label{adp:witness-endpoints}
 r+m\ge t-o(1),\qquad r\le t+o(1).
\end{equation}
No equality between these lengths is assumed.

The inverse moment \cite[Lemma 17.1]{OAI26} has capacity
$z_I(r)=(1-r)/2$ for $r<1$. Applying Theorem~\ref{adp:plain}
to two copies of the plain witness gives
$z_P(m)=(1-2m)/(6\kappa)$ for $m<1/2$.
Order the positive prime amplitudes decreasingly and fill a requested
capacity fractionally in that order. Its gain is at least
$q_{\rm amp}z$. Removing the last fractional factor costs at most
$\delta$ times the largest mesh length in the squared spike.
Approach each positive capacity from below by a fixed decrement.
Both losses may be made arbitrarily small, with the mesh chosen
before the number of physical prime factors.

The witness spikes and the two moments yield the count exponents
\[
 1-\delta r-2q_{\rm amp}z_I(r),\qquad
 1-2\delta m-2q_{\rm amp}z_P(m).
\]
The second denominator is the fourth power of the plain witness.
It is not a second-moment estimate. Substitution, and then
\eqref{adp:witness-endpoints}, give the decreasing and increasing
affine functions
\[
 A_I(r)=1-\delta\{x+(1-x)r\},\qquad
 S_t(r)=1-\delta\{cx+B_x(t-r)\}.
\]
The replacement of $m$ is legitimate because its exponent has
derivative $-2\delta+2\delta cx<0$.
They meet at
\begin{equation}\label{adp:intersection}
 r_*(t)=\frac{B_xt+(c-1)x}{D_x}.
\end{equation}
Their common value is $R_{\rm short}(t)$.
Use the plain count for $r\le r_*$ and the inverse count for
$r\ge r_*$, except in the small zero-capacity neighborhoods.

In the whole range $4/9\le c\le6/13$, $0\le x\le1/2$,
directly from \eqref{adp:intersection},
\[
 r_*(1)\ge23/37,\qquad r_*(3/2)=1,\qquad
 \frac13\le t-r_*(t)\le\frac12 .
\]
Thus the requested inverse length is at most $7/37$. Its second
strict condition is also retained: after decreasing $z_I$ by
$\nu_0>0$,
\[
 1-r-2z\ge2\nu_0,\qquad
 3-2r-8z=4(1-r-2z)+(2r-1)>0.
\]
For the inverse side, $r\ge23/37$ makes the latter margin at least
$9/37$ before its first positive term.
On the plain side $m\ge1/3-o(1)$, so its requested length is at
most $c/6+o(1)\le1/13+o(1)$.
The stated supply margin therefore suffices for both selections.

For $r\ge1$, use the unmarked sixth-power amplification of
\cite[Lemma 17.6, (17.88)]{OAI26}. Its normalized inverse energy
has exponent $e(r)=\max\{1,(1+5r)/6\}$, with arbitrarily small
uniform loss on bounded length ranges. Dividing by the witness
spike gives $1-\alpha+(\alpha-\delta)r$.
Since $\delta\le\kappa<\alpha$, the second inequality in
\eqref{adp:witness-endpoints} bounds this by $L(t)+o(1)$.
For $m\ge1/2$, the zero-factor fourth moment gives exponent
$1-2\delta m\le1-\delta$. This is below $R_{\rm short}(t)$.
If a positive capacity is too small for its decrement, use these
corresponding unweighted bounds. Their uniform forms cover the
shrinking neighborhoods of $r=1$ and $m=1/2$; no marked estimate
is invoked at a vanishing strict margin.

It remains to verify the coefficient and parameter hypotheses.
With $\Theta=\langle\eta,\widehat{\mathcal T}\rangle$, every
physical prime coefficient is a fixed combination
\[
 \nu(p)\one_{p\in1_{\mathcal T}}
   =|\mathcal T|^{-1}
       \sum_{\vartheta\in\widehat{\mathcal T}}(\nu\vartheta)(p),
 \qquad \nu\in\Theta .
\]
If necessary, conjugate the \emph{whole} witness and selected-prime
product; conjugate profiles and zeros with it. This gives the
orientation $\nu(n)\chi_n(u)$ of Theorem~\ref{adp:plain} without
altering an underlying prime support. A retained large sixth-power-free
row has a good prime of valuation $1,\ldots,5$, so its inducing
character ramifies there and lies outside $\Theta$.
Rows supported entirely on $\mathcal S$, including all unit rows,
form the fixed finite set handled by the separate small-row argument.

The selected real parts lie in a bounded smooth family. Recover the
rowwise selected real parts, heights and scales using the common
parameter-Sobolev measure of \cite[Lemma 4.5 and Section 19.2]{OAI26}.
Derivatives insert logarithmic profile weights, and the height
range contributes the fixed power of $1+T_1$.
These selections do not authorize row-dependent prime coefficients.
Fix all internal moment orders before the external recovery-tail
order. This retains the original cutoffs, natural zeros, independent
prime windows and cumulative frequency losses, proving
\eqref{adp:count}.
\end{proof}

For reference, the balance of the two branches is explicit. Put
\[
 \mathcal J=(\alpha-\delta)D_x+\delta P_x>0 .
\]
Then $D_x(R_{\rm short}-L)=\delta P_x/2-\mathcal J(t-1)$, so
\[
 t_*=1+\frac{\delta P_x}{2\mathcal J},\qquad
 R_*=1-\delta+
       \frac{(\alpha-\delta)\delta P_x}{2\mathcal J}.
\]
No replacement by a fixed baseline capacity, and consequently no
additional baseline-capacity loss, is needed. The endpoint
certificate, Euler factors, residues and continuation using this
count are separate from the present supplement.

The remaining elementary branches are unchanged. With no selected
primes, take $t=1$; the unweighted inverse and fourth moments give
count exponent $1-2\delta/3$ by balancing $r=2/3$, $m=1/3$.
The detector's floor bin $a=51/100$ uses the original trivial row
count, of exponent one, rather than a positive-slot estimate.
The fixed finite small-row set is retained separately as described
above. None of these branches needs a new admission or coefficient.

\begin{thebibliography}{9}
\bibitem{OAI26}
\raggedright
OpenAI,
\emph{The Quasi-Riemann Hypothesis: A Zero-Free Half-Plane
$\Re s>7/8$}, unrefereed preprint, 30 September 2026.
All numbered references in this supplement use commit
\texttt{adc7f1241b42e322a6451854ab7e4b4c146bf78a}.
\href{https://github.com/openai/math/blob/adc7f1241b42e322a6451854ab7e4b4c146bf78a/preprints/The-Quasi-Riemann-Hypothesis-September-30-2026/paper.pdf}
{Pinned manuscript}.
The corresponding TeX source has SHA256
\texttt{42a5ee0febca59fd1def55cfd6c6808c32}\allowbreak
\texttt{2ef4237d7726303ea52f711deac6a3}.
\end{thebibliography}
\end{document}
